% !TEX program = xelatex
% UTF-8; standalone source, with no external figures or bibliography.
% Compile twice: xelatex -interaction=nonstopmode -halt-on-error PJ1报告.tex
% Evidence: FDU-AI base commit 385c7d5 plus Q2 best-cost revision, Python 3.12.14, 2026-10-09.
% Command: python autograder.py --no-graphics
% Q9 scores: 1182, 1025, 946, 1102, 867, 1158, 1099, 1043, 1263, 1065.
% Q9 average Hero decision time: 0.763 seconds per game, as printed by the grader.
\documentclass[UTF8,fontset=fandol,a4paper,11pt]{ctexart}
\usepackage[left=23mm,right=23mm,top=20mm,bottom=21mm]{geometry}
\usepackage{amsmath,amssymb,booktabs,tabularx,array}
\usepackage[hidelinks]{hyperref}
\hypersetup{pdftitle={人工智能 PJ1：搜索实验报告},pdfauthor={}}
\setlength{\parindent}{2em}
\setlength{\parskip}{2pt}
\setlength{\emergencystretch}{2em}
\linespread{1.08}
\setlength{\abovedisplayskip}{6pt plus 2pt minus 2pt}
\setlength{\belowdisplayskip}{6pt plus 2pt minus 2pt}
\setlength{\tabcolsep}{6pt}
\ctexset{
  section={format=\large\bfseries,beforeskip=12pt,afterskip=5pt},
  subsection={format=\normalsize\bfseries,beforeskip=8pt,afterskip=3pt}
}
\pagestyle{plain}
\newcommand{\code}[1]{\texttt{#1}}
\newcommand{\md}{d_{\mathrm{M}}}
\newcommand{\mazed}{d_{\mathrm{maze}}}
\newcolumntype{Y}{>{\raggedright\arraybackslash}X}
\begin{document}

\begin{center}
  {\LARGE\bfseries 人工智能 PJ1：搜索实验报告}\par
  \vspace{4pt}
\end{center}
\vspace{-2pt}
\begin{center}
  \begin{tabularx}{0.95\linewidth}{>{\centering\arraybackslash}p{28mm} >{\centering\arraybackslash}p{38mm} Y}
    \toprule
    姓名 & 学号 & 分工 \\
    \midrule
    \rule{0pt}{2mm} 陈家龙 & 24300980041 & Q1、Q2、Q4、Q5、Q6、Q9 \\
    \rule{0pt}{7mm} 刘俊彦 & 24300980021 & Q3、Q7、Q8、撰写报告 \\
    \bottomrule
  \end{tabularx}
\end{center}

\section{实验环境与课上核心算法}

实验运行环境为 Windows、Python 3.12.14。测试全部完成，使用无图形界面的自动评分器。

核心实现位于 \code{search.py}、\code{searchAgents.py} 和 \code{multiAgents.py}。

\subsection{Q1：广度优先搜索}

\code{breadthFirstSearch} 使用 \code{util.Queue} 保存“状态、到达该状态的动作列表”。每次取出队首，先检查目标；首次扩展该状态时加入已访问集合 \code{visited} ，再按原有顺序加入合法后继及其动作路径。队列中可能存在重复状态，但同一状态不会重复扩展。队列按路径深度递增弹出，因此在本实验单位步长代价的迷宫中，首次找到的目标路径步数最少；搜索失败时返回空列表。

\subsection{Q2：A* 搜索}

\code{aStarSearch} 使用 \code{util.PriorityQueue}，以
\[
  f(s)=g(s)+h(s)
\]
为优先级。队列元素保存状态、动作路径和累计路径代价 $g$，用 \code{best\_cost} 记录各状态已发现的最小 $g$ 值。后继只有在新路径代价严格更低时才更新记录并入队；即使状态曾被扩展，也允许通过更短路径重新扩展。弹出时先跳过代价高于当前记录的旧条目，再检查目标并生成后继。因此，在启发函数可采纳且目标值为 0 时，不要求一致性也能保持最优性。$h\equiv0$ 时，该实现等价于一致代价搜索；位置搜索使用曼哈顿距离启发式。

\subsection{Q3：极小极大搜索}

\code{MinimaxAgent} 按智能体索引依次递归：勇士为 0 号，取后继价值最大值；每条恶龙各占一个最小化层，取后继价值最小值。一层搜索深度表示勇士与全部恶龙各行动一次，只在索引循环回勇士时增加深度。胜利、失败或达到指定深度时，调用 \code{self.evaluationFunction} 评价状态。根节点选择价值最高的合法动作，价值相同时保留先遇到的动作。

\section{课后实验实现}

\subsection{Q4：四角问题建模}

\code{CornersProblem} 将状态表示为
\[
  s=\bigl(p,(b_1,b_2,b_3,b_4)\bigr),\qquad b_i\in\{\mathrm{False},\mathrm{True}\},
\]
其中 $p$ 是勇士位置，$b_i$ 表示第 $i$ 个角落是否已访问。起点若处于角落，对应标记初始化为真；四个标记全部为真时达到目标。后继只包含不穿墙的上下左右移动，步长代价为 1，进入角落时更新对应标记。标记使用不可变元组，各搜索状态相互独立，不包含恶龙等无关信息。使用课上完成的 BFS 可直接搜索这个状态空间，得到经过全部角落的最短路径。

\newpage
\subsection{Q5：四角启发函数}

Q5、Q6 均按作业要求构造非负启发函数，并在目标状态返回 0。

设 $R$ 为未访问角落集合，$\md$ 为曼哈顿距离。\code{cornersHeuristic} 枚举 $R$ 的全部排列，取忽略墙壁时访问剩余角落的最短距离和：
\[
  h_C(p,R)=
  \begin{cases}
    0, & R=\varnothing,\\[2pt]
    \displaystyle\min_{\pi\in\mathrm{Perm}(R)}
    \left[\md(p,\pi_1)+\sum_{i=1}^{|R|-1}\md(\pi_i,\pi_{i+1})\right],
    & R\ne\varnothing.
  \end{cases}
\]
最多枚举 $4!=24$ 种顺序。任意真实完成路径都对应一种角落访问顺序，且每段真实距离不小于曼哈顿距离，因此 $h_C$ 不高估剩余代价。对一步后继 $s'$，把该步接在 $s'$ 的最优排列之前；若该步到达未访角落，则先访问该角落。这构成 $s$ 的一个候选方案，故 $h_C(s)\leq1+h_C(s')$。因此该启发式可采纳且一致。

\subsection{Q6：全金币启发函数}

\code{CoinSearchProblem} 的状态由勇士位置 $p$ 与剩余金币网格组成，后继复制网格并清除进入位置的金币。设 $F$ 为剩余金币集合，\code{coinHeuristic} 取到最远金币的真实迷宫距离：
\[
  h_F(p,F)=
  \begin{cases}
    0, & F=\varnothing,\\
    \displaystyle\max_{c\in F}\mazed(p,c), & F\ne\varnothing.
  \end{cases}
\]
\code{mazeDistance} 通过 BFS 求距离，以排序后的两个端点为键，对称缓存于对象的 \code{heuristicInfo} 中；已计算距离在后续查询中可直接复用。

收集全部金币必然到达最远金币，故 $h_F$ 是真实剩余代价的下界。对未被吃掉的金币 $c$，最短路径距离满足 $\mazed(p,c)\leq1+\mazed(p',c)$；若本步吃掉金币，其原距离为 1，同样不超过 $1+h_F(s')$。对各金币取最大值即可得到 $h_F(s)\leq1+h_F(s')$，因此该启发式可采纳且一致。

\subsection{Q7：Alpha-Beta 剪枝}

\code{AlphaBetaAgent} 沿用 Minimax 的轮次与深度定义，并传递祖先的 $\alpha$、$\beta$ 边界。最大化节点随已搜索子节点更新下界，当前价值严格大于 $\beta$ 时剪枝；最小化节点更新上界，当前价值严格小于 $\alpha$ 时剪枝。按 \code{getLegalActions} 的原顺序逐个生成后继，不预先生成全部子节点，也不在相等时剪枝，以符合评分器对节点生成数量的要求。剪枝跳过不影响祖先决策的分支，保留 Minimax 的最优值。

\subsection{Q8：期望极大化}

\code{ExpectimaxAgent} 在勇士节点仍取最大值，在恶龙节点假设每个合法动作等概率发生，使用
\[
  V(s)=\frac{1}{|A(s)|}\sum_{a\in A(s)}V\bigl(\mathrm{Succ}(s,a)\bigr)
\]
代替最小值；多条恶龙按索引依次形成概率层。终止条件和深度计算与 Minimax 相同。Minimax 假设恶龙总选择对勇士最不利的动作，因此更保守；Expectimax 根据随机恶龙的平均结果决策，可能接受存在失败风险但期望收益更高的路线。

\newpage
\subsection{Q9：改进状态评价}

\code{betterEvaluationFunction} 评价当前状态。设 $S$ 为游戏分数，$d_c$、$d_s$ 分别为最近金币、宝剑的曼哈顿距离，$d_i$ 为第 $i$ 条恶龙的曼哈顿距离，$t_i$ 为其剩余虚弱时间。非终局评价为
\[
  E_0(s)=S+C(s)+B(s)+\sum_i\frac{w(t_i)}{d_i+1},
\]
其中，存在金币时 $C(s)=1/(d_c+1)$，否则为 0；存在宝剑时 $B(s)=7/(d_s+1)$，否则为 0。恶龙权重为
\[
  w(t)=
  \begin{cases}
    -10, & t=0,\\
    5, & t>5,\\
    -2, & 0<t\leq5.
  \end{cases}
  \qquad
  E(s)=
  \begin{cases}
    S+100000, & \text{胜利},\\
    S-100000, & \text{失败},\\
    E_0(s), & \text{非终局}.
  \end{cases}
\]
游戏分数提供已取得收益的信号；距离奖励引导接近金币与宝剑；正常恶龙受到较强惩罚，虚弱时间充足时鼓励追击，虚弱即将结束时改为轻度回避。分母加 1 避免距离为 0 时除零，终局的大幅加减分优先区分胜负。评价使用曼哈顿距离以控制计算开销，其距离信号可能忽略墙壁绕行，是当前设计的局限。

\section{自动评分结果}

在仓库根目录、Python 3.12.14 环境下执行：
\begin{quote}
  \small\ttfamily python autograder.py --no-graphics
\end{quote}
现有自动评分用例全部通过，结果如下。表中的“扩展”指对应外层搜索问题的扩展节点数。Q9 使用 \code{smallClassic}、深度为 2 的 \code{ExpectimaxAgent} 和 \code{evalFn=better}，与一条随机恶龙对战；随机种子设为 0，连续运行 10 局。决策时间为平均每局勇士累计思考时间。

\begin{center}
  \small
  \renewcommand{\arraystretch}{1.22}
  \begin{tabularx}{\linewidth}{@{}p{10mm}Y r@{}}
    \toprule
    题目 & 完成情况与主要验证结果 & 得分 \\
    \midrule
    Q1 & BFS 已实现；\code{mediumMaze}：68 步，扩展 269 个节点 & 3/3 \\
    Q2 & A* 已实现；同图使用曼哈顿启发式：68 步，扩展 221 个节点 & 3/3 \\
    Q3 & Minimax 已实现；通过仓库现有评分用例 & 5/5 \\
    Q4 & 四角状态建模完成；\code{tinyCorners} 最短路径为 28 步 & 3/3 \\
    Q5 & 四角启发式完成；\code{mediumCorners}：106 步，扩展 741 个节点 & 3/3 \\
    Q6 & 金币启发式完成；\code{trickySearch} 扩展 4137 个节点 & 4/4 \\
    Q7 & Alpha-Beta 已实现；通过动作、深度及节点生成检查 & 5/5 \\
    Q8 & Expectimax 已实现；通过概率节点、多恶龙及深度检查 & 5/5 \\
    Q9 & 改进评价完成；10/10 局获胜，平均得分 1075.0，平均决策时间 0.763 秒，均未超时 & 6/6 \\
    \midrule
    合计 & 课上 Q1--Q3：11/11；课后 Q4--Q9：26/26 & \textbf{37/37} \\
    \bottomrule
  \end{tabularx}
\end{center}

\end{document}
